\documentclass[11pt]{article}

\usepackage[a4paper,margin=1in]{geometry}
\usepackage{amsmath,amssymb,bm,mathtools}
\usepackage{booktabs}
\usepackage{hyperref}
\usepackage{microtype}

\hypersetup{
  colorlinks=true,
  linkcolor=blue,
  citecolor=blue,
  urlcolor=blue
}

\newcommand{\ii}{\mathrm{i}}
\newcommand{\dd}{\mathrm{d}}
\newcommand{\bfu}{\bm{u}}
\newcommand{\bfq}{\bm{q}}
\newcommand{\bfzero}{\bm{0}}
\newcommand{\Graw}{G_{\mathrm{raw}}}
\newcommand{\Gtilde}{\widetilde{G}}
\newcommand{\Ex}{E_x}
\newcommand{\Ey}{E_y}
\newcommand{\Ez}{E_z}
\newcommand{\Dy}{D_y}
\newcommand{\Bbdry}{B_{\partial}}

\title{Planned numerical sampling and constraint construction for an MKM channel DNS at $Re_{\tau}=395$}
\author{Prepared for the Shenfun MKM channel-flow workflow}
\date{27 July 2026}

\begin{document}

\maketitle

\begin{abstract}
This note specifies a gated plan for a new turbulent-channel direct numerical
simulation at nominal friction Reynolds number $Re_{\tau}=395$.  It adapts the
existing two-stage Shenfun MKM workflow to the mesh and domain used by Moser,
Kim, and Mansour for their corresponding case.  Particular attention is paid
to the reference paper's odd wall-normal count, the distinction between the reference
Chebyshev--Gauss--Lobatto convention and Shenfun's Chebyshev--Gauss quadrature
grid, the even-size requirement of the optimized Shenfun biharmonic solver,
the constant-flux target, integer de-aliasing sizes, measured wall stress,
sampling convergence, constraint construction, and the much larger storage
cost of an all-plane temporal record.  The proposed persistent representation
stores the two independent MKM spectral state variables and the required zero
horizontal modes, from which all three velocities on all 192 horizontal
planes are reconstructed without approximation beyond solver roundoff.  This
is a planning document: no $Re_{\tau}=395$ production simulation or result is
claimed here.
\end{abstract}

\noindent\textbf{Status:} the Ubuntu and CentOS 18-plane and 192-plane
numerical/I/O gates are verified.  The former full-resolution
$\Delta t=2.5\times10^{-4}$ pilot was unstable; the refined-step pilot,
physical stationarity, and resolution-acceptance gates remain pending.

\section{Objective and reference case}

Let $x\in[0,L_x)$ be streamwise, $z\in[0,L_z)$ be spanwise, and
$y\in[-h,h]$ be wall-normal.  The proposed nondimensional channel has
\begin{equation}
  h=1,\qquad L_x=2\pi h,\qquad L_z=\pi h,
  \label{eq:domain}
\end{equation}
with periodicity in $x$ and $z$ and no slip at $y=\pm h$.  Velocities are
ordered as
\[
  \bfu=(u_x,u_z,u_y)^T
\]
in the target files: streamwise, spanwise, and wall-normal.  The internal MKM
HDF5 groups remain
\[
  \texttt{u0}=u_y,\qquad \texttt{u1}=u_x,\qquad
  \texttt{u2}=u_z.
\]

Moser, Kim, and Mansour reported the parameters in
Table~\ref{tab:mkm_reference}.  Their nominal case name is $Re_{\tau}=395$,
while the measured value in the paper and statistical database is 392.24.

\begin{table}[ht]
  \centering
  \caption{Moser--Kim--Mansour reference parameters for the nominal
  $Re_{\tau}=395$ channel.}
  \label{tab:mkm_reference}
  \begin{tabular}{lc}
    \toprule
    Quantity & Reference value \\
    \midrule
    Actual $Re_{\tau}$ & 392.24 \\
    $(L_x,L_z)/h$ & $(2\pi,\pi)$ \\
    Paper order $N_x\times N_y\times N_z$ & $256\times193\times192$ \\
    $\Delta x^+$ & approximately 10.0 \\
    $\Delta z^+$ & approximately 6.5 \\
    $\Delta y_c^+$ & approximately 6.5 \\
    \bottomrule
  \end{tabular}
\end{table}

The source paper states that domain sizes were chosen so two-point
correlations were essentially zero at half-box separation and that resolution
was chosen so energy spectra were sufficiently small at the largest retained
wavenumbers.  It also reports at least 13 wall-normal points below $y^+=10$.
These are validation conditions, not merely historical metadata.

The primary reference data are available from the University of Texas
turbulence database:
\begin{center}
  \url{https://turbulence.oden.utexas.edu/MKM_1999.html}.
\end{center}

\section{Proposed Shenfun mesh}

\subsection{Ordering and nominal spacings}

The current driver uses solver order
\[
  (N_{\mathrm{wall}},N_{\mathrm{stream}},N_{\mathrm{span}}).
\]
The user-selected Shenfun baseline is
\begin{equation}
  \boxed{\texttt{--n 192 256 192}}.
  \label{eq:solver_n}
\end{equation}
At nominal $Re_{\tau}=395$, its two uniform-direction spacings are
\begin{align}
  \Delta x^+
  &=\frac{L_x}{N_{\mathrm{stream}}}\frac{u_{\tau}}{\nu}
    =\frac{2\pi(395)}{256}=9.69476,\\
  \Delta z^+
  &=\frac{L_z}{N_{\mathrm{span}}}\frac{u_{\tau}}{\nu}
    =\frac{\pi(395)}{192}=6.46317.
  \label{eq:horizontal_spacing}
\end{align}

The saved wall-normal mesh is Shenfun's Chebyshev--Gauss quadrature grid,
\begin{equation}
  y_m=\cos\left(\frac{(2m+1)\pi}{2N_{\mathrm{wall}}}\right),
  \qquad m=0,\ldots,N_{\mathrm{wall}}-1.
  \label{eq:gauss_nodes}
\end{equation}
For even $N_{\mathrm{wall}}=192$, the centerline lies midway between the two
central Gauss points.  The gap across the center, distance of either central
point from the center, and first saved distance from a wall are
\begin{align}
  \Delta y_c^+
  &=2(395)\sin\left(\frac{\pi}{2(192)}\right)=6.46310,\\
  y_{\mathrm{nearest\ center}}^+
  &=395\sin\left(\frac{\pi}{2(192)}\right)=3.23155,\\
  y_{1}^{+}
  &=395\left[1-\cos\left(\frac{\pi}{2(192)}\right)\right]
    =0.01322.
  \label{eq:wall_spacing}
\end{align}
There are 14 positive-half-channel Gauss nodes with distance below
$y^+=10$ from the nearest wall.

\subsection{Why the paper uses 193 but Shenfun uses 192}

The paper does not identify an odd count as a separate physical requirement.
The most direct numerical explanation is the conventional
Chebyshev--Gauss--Lobatto representation.  If the maximum polynomial degree is
even, the number of endpoint-including nodal values is one greater and hence
odd.  Counts 129, 193, and 257 correspond to even interval/degree counts 128,
192, and 256; this also includes both walls and the centerline and is
convenient for cosine-transform implementations.

Shenfun's grid in Eq.~\eqref{eq:gauss_nodes} is not the same nodal family:
the wall values are excluded and no slip is imposed through composite
Galerkin bases.  Although its transforms and bases can represent odd sizes,
the optimized Chebyshev biharmonic solver used by the MKM
wall-normal-velocity equation splits the coefficient system into equal
even/odd parity blocks.  A four-rank Ubuntu test at
$N_{\mathrm{wall}}=17$ was finite after initialization but generated
non-finite velocity coefficients in the first Runge--Kutta substage.  Its
matched $N_{\mathrm{wall}}=18$ control remained finite for 40 steps and
reconstructed every saved velocity value exactly from the independent state.

The current solver must therefore use an even wall-normal size.  The selected
$N_{\mathrm{wall}}=192$ is one point, or approximately 0.52 percent, below
the paper's count.  It retains 14 saved points per wall below $y^+=10$ and
gives a center gap essentially equal to the reference center spacing.

The two Fourier counts should remain even.  This retains efficient FFT sizes,
the expected real-transform and Nyquist conventions, and convenient MPI
factorizations.

\subsection{Integer de-aliasing size}

The nonlinear terms use padded transforms.  The current implementation
replaces a requested padding factor $p$ by an integer size
\[
  N_p=\lfloor pN\rfloor.
\]
A literal $p=1.5$ at $N=192$ gives the exact integer padded size 288.  The proposed
integer padded sizes are
\begin{equation}
  (N_{p,y},N_{p,x},N_{p,z})=(288,384,288),
  \label{eq:padded_sizes}
\end{equation}
so the run card should use
\begin{equation}
  \texttt{--padding-factor 1.5 1.5 1.5}.
  \label{eq:padding_factors}
\end{equation}

\section{Reynolds number and constant-flux target}

The nominal nondimensional inputs are
\begin{equation}
  h=1,\qquad u_{\tau,\mathrm{ref}}=1,\qquad
  \nu=\frac{1}{395}.
  \label{eq:viscosity}
\end{equation}
However, prescribing $\nu$ does not by itself demonstrate that the developed
constant-flux solution has $Re_{\tau}=395$.  If the time-averaged wall shear is
$\overline{\tau}_w$, the measured friction velocity and Reynolds number are
\begin{equation}
  u_{\tau,\mathrm{meas}}=\sqrt{\overline{\tau}_w},\qquad
  Re_{\tau,\mathrm{meas}}
  =\frac{u_{\tau,\mathrm{meas}}h}{\nu}.
  \label{eq:measured_retau}
\end{equation}
Both walls must be evaluated and their symmetry reported.

The current $Re_{\tau}=180$ runner defaults to an $Re_{\tau}=180$ bulk velocity and must
not be reused unchanged.  Numerical integration of the official MKM
$Re_{\tau}=392.24$ mean profile gives the planning value
\begin{equation}
  \frac{U_b}{u_{\tau}}=17.5447511436.
  \label{eq:bulk_velocity}
\end{equation}
For Eq.~\eqref{eq:domain}, the corresponding integrated flux is
\begin{equation}
  Q=2hL_xL_zU_b=4\pi^2 U_b=692.639012411.
  \label{eq:integrated_flux}
\end{equation}
This is an initial calibration target.  The developed simulation is accepted
only if Eq.~\eqref{eq:measured_retau} is within the agreed tolerance; otherwise
$U_b$ is adjusted and a calibration continuation is run.

\section{Initialization changes required before execution}

Changing only \texttt{--re} is insufficient.  The present
\path{demo/MKM.py} contains an $Re_{\tau}=180$ initial centerline/profile constant,
the old bulk-velocity default, and perturbation choices whose wall-unit
wavelengths change with Reynolds number.  Before production, the run metadata
must record:
\begin{enumerate}
  \item nominal $Re_{\tau}$, $\nu$, target $U_b$, and integrated flux;
  \item the initial mean profile and its source;
  \item perturbation amplitudes and wavelengths in wall or outer units;
  \item the realized integer padded sizes;
  \item a checksum or version identifier for the reference profile.
  \item an all-plane state output path for the independent wall-normal
        velocity and wall-normal vorticity coefficients, the two separately
        evolved zero-horizontal-mode profiles, online block moments, and
        infrequent complete audit fields.  The current runner exposes only one
        full-field \texttt{modsave} cadence, so this recorder and its decoder
        must be implemented and smoke-tested before stationary sampling.
\end{enumerate}

The recommended clean start interpolates the official MKM mean profile onto
Eq.~\eqref{eq:gauss_nodes}, reflects it across the centerline, and adds small
resolved perturbations.  The interpolated profile is only an initial
condition and validation reference.  It is not counted as new DNS evidence.
The initialized velocity must pass divergence and no-slip checks after the
same forward/backward projections used by the solver.

\section{Preflight gates}
\label{sec:gates}

\subsection{Even-grid and restart smoke}

First run a small even wall-normal case
\[
  (N_{\mathrm{wall}},N_{\mathrm{stream}},N_{\mathrm{span}})
  =(18,32,24),
\]
through initialization, a forced checkpoint, restart, quadrature HDF5 output,
state reconstruction, and a complete NaN/Inf audit.  The uninterrupted and
restarted states must agree to solver tolerance.

\subsection{Full-shape decomposition benchmark}

Allocate the unpadded shape $(192,256,192)$ and padded shape
$(288,384,288)$ on the Linux server and take only a few steps without
snapshots.  Benchmark feasible MPI sizes, initially 16, 32, and 64 ranks.
The production rank count is selected from successful decomposition, memory,
and step-time measurements.

The isolated Ubuntu pilot completed two full-shape steps at all three rank
counts.  The 64-rank one-sample I/O gate then reconstructed all 192 planes
with maximum absolute component error $2.14\times10^{-14}$, maximum relative
error $1.06\times10^{-15}$, maximum divergence $9.53\times10^{-16}$, and no
non-finite values.  The independent-state and physical-velocity files were
152,594,328 and 226,523,696 bytes, respectively.  A reduced spin-up followed
by two sampling segments also passed checkpoint restart, HDF5 append,
state-shard continuity, and reconstruction checks.  CentOS jobs 3201--3204
subsequently passed the 16-, 32-, and 64-rank allocation gates and the
64-rank state-I/O gate.  The requested scheduler policy allocates two SLURM
tasks per launched MPI rank: a 64-rank run requests
\texttt{\#SBATCH --ntasks=128}, while the launcher explicitly uses
\texttt{mpiexec -n 64}.  The shared script validates and records this 2:1
allocation-to-launch ratio.

\subsection{Time-step pilot}

The full-resolution CentOS pilot in job 3206 used
\begin{equation}
  \Delta t=2.5\times10^{-4}.
  \label{eq:initial_dt}
\end{equation}
This step is rejected.  The reported CFL exceeded one at $t=0.1$, approached
1.6, and the spanwise kinetic energy grew from 0.3907 to 12.8004 by $t=0.8$.
At $t=0.9$, all 28,311,552 physical velocity values were non-finite.  The
minimum Chebyshev--Gauss wall spacing at $N_{\mathrm{wall}}=192$ is about one
ninth of that at 64 points.  Scaling the old $Re_{\tau}=180$ step
$5\times10^{-4}$ by this ratio gives $5.56\times10^{-5}$.  The conservative
candidate used in the revised scripts is therefore
\begin{equation}
  \boxed{\Delta t=5\times10^{-5}}.
  \label{eq:refined_dt}
\end{equation}
The first full-mesh refined pilot ends at $t=1$ and records diagnostics every
0.1 outer units.  Finite values, CFL margin, energy, divergence, restart
consistency, and high-wavenumber spectra determine acceptance.

The refined step and cadence passed a four-rank, 200-step
$(18,32,24)$ Ubuntu round trip: all 20 retained samples reconstructed exactly,
the finite-value audit passed, maximum divergence was
$1.06\times10^{-15}$, and CFL remained at or below 0.00609.  This validates
the revised software and I/O cadence, not full-mesh stability.

\subsection{Spatial-resolution acceptance}

The full-grid pilot must demonstrate:
\begin{enumerate}
  \item no high-wavenumber energy pile-up in either periodic direction;
  \item mean and Reynolds-stress profiles consistent with the MKM database;
  \item streamwise and spanwise two-point correlations nearly zero at
        half-box separation;
  \item at least 13 points below $y^+=10$;
  \item symmetry of the two channel halves;
  \item small divergence and no-slip constraint residuals.
\end{enumerate}

\section{Two-stage production sampling}

\subsection{Spin-up}

The development stage writes no target snapshots and forces a checkpoint at
the end of each restartable segment.  It records bulk flux, wall shear,
measured $Re_{\tau}$, all three kinetic-energy components, divergence, CFL,
mean profiles, Reynolds stresses, and spectral-tail ratios.  After the
refined-step $t=1$ calibration, use five-outer-unit restart segments when the
measured wall time fits the allocation.  Candidate physical reviews are made
at outer times 20, 40, 60, 80, 100, and 120.  The user-selected
sampling start is fixed at $t=120$, so no retained temporal sampling begins
earlier.  The old $t=20$ spin-up is not automatically transferred to the
higher-Reynolds-number case.

A checkpoint is accepted only when bulk flux and measured $Re_{\tau}$ are
steady, both walls agree, energy components show no block drift, mean and
Reynolds profiles have settled, and spectral/correlation checks pass.  A
planning calibration criterion is
\begin{equation}
  \frac{|Re_{\tau,\mathrm{meas}}-395|}{395}\le0.02.
  \label{eq:retau_tolerance}
\end{equation}

\subsection{Stationary sampling}

The sampling stage restarts from the accepted $t=120$ checkpoint into a
separate velocity file and ends at $t=300$.  Thus
\begin{equation}
  t_{\mathrm{sample,start}}=120,\qquad
  t_{\mathrm{sample,end}}=300,\qquad
  T_{\mathrm{sample}}=180.
  \label{eq:sampling_window}
\end{equation}
At the refined step these endpoints are DNS steps 2,400,000 and 6,000,000,
respectively, so the stationary stage contains 3,600,000 DNS steps.
The dense cadence used for temporal observables is
\begin{equation}
  \boxed{\Delta t_{\mathrm{temporal}}u_{\tau}/h=0.005}.
  \label{eq:sampling_cadence}
\end{equation}
For Eq.~\eqref{eq:refined_dt}, this is one temporal sample every 100 DNS steps
and
\begin{equation}
  \Delta t_{\mathrm{temporal}}^+
  =Re_{\tau}\Delta t_{\mathrm{temporal}}=1.975.
  \label{eq:sampling_wall_time}
\end{equation}
The corresponding temporal Nyquist angular frequency is
\begin{equation}
  \omega_{\mathrm{Ny}}h/u_{\tau}
  =\frac{\pi}{\Delta t_{\mathrm{temporal}}}=628.319.
  \label{eq:sampling_nyquist}
\end{equation}
Dense full fields are not the primary temporal record.  At the dense cadence
the writer stores the complete independent MKM spectral state, from which all
three velocity components are reconstructed on every one of the 192
wall-normal planes during postprocessing.  All-wall mean and Reynolds-stress
raw moments are also accumulated by time block as an independent audit.  A
complete physical audit velocity field is retained only every 10 outer time
units, or every 200000 DNS steps, and at the final accepted state.  State data
are written in 1000-sample shards (five outer time units), with a forced
checkpoint every five outer units.  Convergence is reviewed over accumulated
40-outer-time-unit blocks.

A short pilot records temporal observables every 50 DNS steps
($\Delta t=0.0025$, $\Delta t^+=0.9875$) and compares them with the same
record decimated to 0.005.  The production cadence is reduced to 0.0025 only
if the 0.005 correlations or spectra show meaningful aliasing in the frequency
band of interest.

For the fixed 180-unit stationary window,
Eq.~\eqref{eq:sampling_cadence} gives 36,000 samples.  The
all-plane temporal auto-correlations and auto-spectra are evaluated from
reconstructed float64 velocity batches in the same time order used by the
existing estimators.  A Hann/Welch estimate with 8,192-sample segments and 50
percent overlap has segment duration 40.96 and angular-frequency spacing about
0.1534.  A full-record diagnostic periodogram has spacing about 0.03491.

The datasets \path{lag_covariance/lag_0} and
\path{lag_covariance/lag_1} are not required and will not be generated.  Any
legacy postprocessor invocation uses \texttt{--max-lag -1}.

Candidate retained windows are evaluated with at least 12 blocks.  The first
field-statistics gate is the established five-percent tolerance for both mean
and Reynolds-stress profiles.  A three-percent sensitivity check is preferred
before final acceptance.  Scalar energy drift is reported separately rather
than hidden by the field-profile result.

\subsection{Exact independent-state record}

The KMM implementation advances the wall-normal velocity coefficient field
$\widehat{u}_y$ in its biharmonic space and the wall-normal vorticity
$\widehat{\eta}_y$ in its Dirichlet space.  For a nonzero horizontal
wavenumber pair $(\kappa,\lambda)$, the same reconstruction used by
\path{demo/ChannelFlow.py} is
\begin{align}
  \widehat{u}_x
  &=\ii\frac{\kappa D_y\widehat{u}_y+
                    \lambda\widehat{\eta}_y}
                   {\kappa^2+\lambda^2},\\
  \widehat{u}_z
  &=\ii\frac{\lambda D_y\widehat{u}_y-
                    \kappa\widehat{\eta}_y}
                   {\kappa^2+\lambda^2}.
  \label{eq:state_reconstruction}
\end{align}
The two $(\kappa,\lambda)=(0,0)$ profiles,
$\widehat{u}_x(y;0,0)$ and $\widehat{u}_z(y;0,0)$, are evolved by separate
one-dimensional equations and must be saved explicitly.  Therefore each
temporal sample contains
\[
  \left\{\widehat{u}_y,\widehat{\eta}_y,
  \widehat{u}_x(\,\cdot\,;0,0),
  \widehat{u}_z(\,\cdot\,;0,0)\right\}.
\]
Restoring these objects and applying the solver's own derivative, wavenumber,
zero-mode insertion, and backward-transform paths recovers every velocity
component at every saved wall-normal point.  The state record uses
complex128/float64 and the native real-to-complex Fourier half-spectrum.
Wall-normal subsampling, horizontal mode truncation, reduced precision, and
lossy compression are excluded from the baseline.

\section{Sampled target objects}

For retained samples $t_n$, the horizontal/time mean is
\begin{equation}
  \overline{u}_{\alpha}(y_m)
  =\frac{1}{N_sN_xN_z}
  \sum_{n,j,\ell}
  u_{\alpha}(x_j,z_{\ell},y_m,t_n).
  \label{eq:mean}
\end{equation}
With $u'_{\alpha}=u_{\alpha}-\overline{u}_{\alpha}$, the Reynolds profile is
\begin{equation}
  R_{\alpha\beta}(y_m)
  =\frac{1}{N_sN_xN_z}
  \sum_{n,j,\ell}
  u'_{\alpha}(x_j,z_{\ell},y_m,t_n)
  u'_{\beta}(x_j,z_{\ell},y_m,t_n).
  \label{eq:reynolds}
\end{equation}

The normalized horizontal Fourier coefficients are
\begin{equation}
  \widehat{u}_{\alpha}(y_m;\kappa_p,\lambda_q,t_n)
  =\frac{1}{N_xN_z}\sum_{j,\ell}
  u'_{\alpha}(x_j,z_{\ell},y_m,t_n)
  e^{-\ii(\kappa_px_j+\lambda_qz_{\ell})}.
  \label{eq:fft}
\end{equation}
For one mode the target stacks the three components level by level into
\begin{equation}
  \widehat{\bfq}_{pq}
  =[\widehat{u}_x(y_0),\widehat{u}_z(y_0),\widehat{u}_y(y_0),\ldots,
    \widehat{u}_x(y_{191}),\widehat{u}_z(y_{191}),
    \widehat{u}_y(y_{191})]^T
  \in\mathbb{C}^{576}.
  \label{eq:modal_vector}
\end{equation}
The equal-time covariance is
\begin{equation}
  B^{(0)}_{pq}
  =\frac{1}{N_s}\sum_n
  \widehat{\bfq}_{pq}(t_n)\widehat{\bfq}_{pq}(t_n)^*.
  \label{eq:b0}
\end{equation}

\section{Constraint on the even Chebyshev--Gauss grid}

For each horizontal mode, define component selectors
\[
  \Ex\widehat{\bfq}=\widehat{\bm{u}}_x,\qquad
  \Ey\widehat{\bfq}=\widehat{\bm{u}}_y,\qquad
  \Ez\widehat{\bfq}=\widehat{\bm{u}}_z,
\]
where here $\Ey$ selects wall-normal velocity and $\Ez$ selects spanwise
velocity in physical notation.  In the target implementation the selector
names must follow the stored component metadata rather than ambiguous
coordinate letters.

Let $\Dy$ be the barycentric nodal derivative on the 192 Gauss points and let
$\Bbdry\in\mathbb{R}^{2\times192}$ interpolate those points to $y=\pm1$.
The modal divergence and boundary rows are
\begin{align}
  G_{\mathrm{div}}(\kappa,\lambda)
  &=\ii\kappa E_{\mathrm{stream}}
    +\ii\lambda E_{\mathrm{span}}
    +\Dy E_{\mathrm{wall}},\\
  G_{\partial}
  &=\begin{bmatrix}
       \Bbdry E_{\mathrm{stream}}\\
       \Bbdry E_{\mathrm{span}}\\
       \Bbdry E_{\mathrm{wall}}
     \end{bmatrix}.
  \label{eq:constraint_blocks}
\end{align}
Thus
\begin{equation}
  \Graw(\kappa,\lambda)
  =\begin{bmatrix}G_{\mathrm{div}}\\G_{\partial}\end{bmatrix}
  \in\mathbb{C}^{198\times576}.
  \label{eq:graw}
\end{equation}
If
\[
  \Graw=U\Sigma V^*,
\]
and $V_r$ contains the right singular vectors retained at the documented
relative tolerance, the saved row-orthonormal constraint is
\begin{equation}
  \Gtilde=V_r^*,\qquad
  \Gtilde\Gtilde^*=I,qquad
  \Gtilde\widehat{\bfq}=0.
  \label{eq:gtilde}
\end{equation}
The expected generic row count is 198, with two dependencies commonly present
at the zero horizontal mode; the actual rank table must be computed and saved,
not assumed.

\section{Storage scaling and processing decision}

The proposed mesh has 49,152 horizontal modes and modal dimension 576.  The
uncompressed estimates in Table~\ref{tab:storage} show that the old output
policy cannot be copied blindly.

\begin{table}[ht]
  \centering
  \caption{Approximate uncompressed storage for the proposed mesh.}
  \label{tab:storage}
  \begin{tabular}{ll}
    \toprule
    Object & Size \\
    \midrule
    One three-component physical snapshot & 0.210938 GiB \\
    36,000 physical velocity samples & 7.4158 TiB \\
    One independent spectral-state sample & 0.142093 GiB \\
    36,000 independent state samples & 4.9954 TiB \\
    18 complete audit fields & 3.797 GiB \\
    Full equal-time \path{modal/B0_DNS} & approximately 243 GiB \\
    \bottomrule
  \end{tabular}
\end{table}

For the real-to-complex horizontal layout, the nonredundant coefficient extent
is $192\times256\times97$.  Two complex128 state arrays and two float64
zero-mode profiles contain 152,570,880 bytes per sample.  This exact
representation saves approximately 32.64 percent relative to three dense
physical components while retaining every horizontal plane.  State files are
split into 1000-sample shards, each approximately 142.09 GiB before
compression.

Only representation-preserving reductions are admitted: the native Fourier
half-spectrum, the independent state variables, and benchmarked lossless
compression.  Compression ratios are not assumed in capacity planning.
Float32, lossy filters, mode truncation, wall-normal subsampling, and further
temporal decimation are excluded unless a separate error study is accepted.
The server allocation should be at least 7 TiB for the baseline state,
checkpoints, audit fields, derived outputs, and safety headroom, excluding any
large disposable reconstruction cache.  A full $B_0$ requires an explicit
scientific need and an additional I/O benchmark.  Lag covariances are omitted
unconditionally.

The detailed sufficient-data construction and its numerical-equivalence test
are recorded in
\path{scripts/mkm_dns_target/TEMPORAL_SAMPLING_REMEDY_Re_tau_395.md}.

Before production, at least 512 samples are dual-written as conventional
physical fields and as independent state.  Reconstructed coefficients,
physical velocities, mean/Reynolds profiles, and temporal correlations and
spectra at all 192 planes must agree to bitwise identity where possible and
otherwise to a documented roundoff/FFT-order tolerance.  The test must include
the zero horizontal mode, two MPI decompositions, shard interruption recovery,
lossless decompression checksums, and measured write throughput.  The existing
constraint builder and postprocessors must also be benchmarked and, if
necessary, optimized at 49,152 horizontal modes.

\section{Planned command card}

After the preflight gates have passed, the common run settings are intended to
resolve as follows:
\begin{verbatim}
TMPDIR=/media/jay/data1/tmp
ENV=/media/jay/data1/conda_envs/shenfun_dns_np126_20260702
REPO=/media/jay/data1/shenfun
OUT=/media/jay/data1/shenfun_dns_runs/
    production_mkm_Re395_N192_256_192
PREFIX=MKM_production_Re395_N192_256_192

--n 192 256 192
--domain -1 1 0 6.283185307179586 0 3.141592653589793
--re 395
--bulk-velocity 17.54475114359453
--dt 0.00005
--padding-factor 1.5 1.5 1.5
--family C
--timestepper IMEXRK222
--save-mesh quadrature
\end{verbatim}
The 64-rank candidate reserves 128 SLURM tasks.  The refined time step remains
subject to the $t=1$ acceptance pilot.  The selected spin-up end and retained
sampling window are fixed at $t=120$ and $120\le t\le300$; the scripts reject
checkpoints or segment endpoints outside that schedule.

\section{Acceptance record}

The final handoff will include the code revision, resolved initialization,
integer padded sizes, server environment, MPI benchmark, selected $\Delta t$,
measured CFL, time-averaged wall stresses and actual $Re_{\tau}$, accepted
spin-up checkpoint, sampling-window convergence, reference-profile and
spectral comparisons, half-box correlations, constraint ranks and residuals,
file sizes/checksums, the exact dense temporal cadence and estimator settings,
the state-reconstruction version and round-trip report, and a list of
deliberately omitted large products.  It will explicitly record that
\path{lag_0} and \path{lag_1} were not generated.

\section{Conclusion}

The reference paper's odd count is a Lobatto-grid convention, but the current
optimized Shenfun MKM biharmonic solver requires an even wall-normal size.
The selected $N_{\mathrm{wall}}=192$ preserves the reference center resolution
closely and is vastly finer than the old 64-point mesh, while Shenfun continues
to impose the walls through its Galerkin bases.  The decisive remaining
prerequisites are an
$Re_{\tau}=395$-aware initialization and flux target, explicit upward rounding of
the de-aliasing sizes, measured wall-stress diagnostics, the refined
full-shape time-step pilot, and a storage-aware postprocessing policy.  Long
production should begin only after the refined-step and physical-acceptance
gates have passed.

\end{document}
